% 星座（综述）
% license CCBYSA3
% type Wiki

本文根据 CC-BY-SA 协议转载翻译自维基百科\href{https://en.wikipedia.org/wiki/Constellation}{相关文章}。

\begin{figure}[ht]
\centering
\includegraphics[width=6cm]{./figures/a943f80cad7e3c1f.png}
\caption{} \label{fig_Coti_1}
\end{figure}
星座是天球上的一个区域，在该区域内，一组可见的恒星形成了人们所感知的图案或轮廓，通常代表一种动物、神话题材，或无生命的物体。$^{[1]}$

最早的星座很可能是在史前时期被定义的。人们利用星座来讲述关于其信仰、经历、创世观念以及神话的故事。不同的文化和国家创造了各自的星座体系，其中一些一直延续到20世纪初，直到当今的星座体系在国际上获得统一认可。随着时间的推移，对星座的认知发生了显著变化。许多星座的大小或形状发生了改变。有些一度流行，后来却逐渐被遗忘。有些则仅局限于单一文化或国家。对星座进行命名也帮助天文学家和航海者更容易地识别恒星。$^{[2]}$

十二个（或十三个）古代星座属于黄道星座，它们横跨黄道——太阳、月球和行星运行所经过的路径。黄道的起源在历史上仍不确定；其占星学上的划分大约在公元前400年于巴比伦或迦勒底天文学中变得突出。$^{[3]}$ 星座通过希腊文化进入西方世界，并出现在赫西俄德、欧多克索斯和阿拉托斯的著作中。由黄道星座及另外36个星座组成的传统48星座（由于阿尔戈船座被划分为三个星座，如今为38个），由出生于埃及亚历山大的希腊—罗马天文学家托勒密在其著作《天文学大成》中列出。星座的形成是大量神话的主题，其中最著名的当属拉丁诗人奥维德的《变形记》。15世纪至18世纪中叶，随着欧洲探险者开始前往南半球，远南天区的星座被陆续加入。由于罗马和欧洲文化的传播，每个星座都拥有一个拉丁名称。

1922年，国际天文学联合会正式接受了现代的88个星座列表，并于1928年采用了官方的星座边界，这些边界共同覆盖了整个天球。$^{[4][5]}$ 在天球坐标系统中，任意一个给定点都位于某一个现代星座之内。一些天文学命名体系在命名天体时会包含其所在的星座，以传达该天体在天空中的大致位置。例如，恒星的弗兰斯蒂德命名法由一个数字和星座名称的所有格形式组成。

其他被称为星群的恒星图案或组合，在严格定义下并不属于星座，但同样被观测者用来辨认夜空。星群可能是某一星座中的几颗恒星，也可能与多个星座共享恒星。星群的例子包括人马座中的茶壶，以及大熊座中的北斗七星。$^{[6][7]}$
\subsection{术语}  
“星座”一词源自晚期拉丁语 cōnstellātiō，可译为“恒星的集合”；该词于14世纪进入中古英语。$^{[8]}$ 古希腊语中表示星座的词是 ἄστρον（astron）。这些术语在历史上泛指任何可被识别的恒星图案，其外观通常与神话人物或生物、地面动物或某种物体相关联。$^{[1]}$ 随着时间推移，在欧洲天文学家之间，星座逐渐被明确界定并获得广泛认可。20世纪，国际天文学联合会（IAU）正式确认了88个星座。$^{[9]}$

当从地球上某一特定纬度观测时，若某个星座或恒星永远不会落到地平线以下，则称其为拱极星座。从北极或南极观测，位于天赤道以南或以北的所有星座都是拱极的。根据不同定义，赤道星座可以指位于北纬45°至南纬45°之间的星座$^{[10]}$，也可以指穿过黄道（或黄道带）赤纬范围、即北纬23.5°至南纬23.5°之间的星座。$^{[11][12]}$

星座中的恒星在天空中看似彼此接近，但它们通常位于距离地球各不相同的位置。由于每颗恒星都具有各自独立的运动，所有星座都会随着时间缓慢发生变化。经过数万至数十万年，熟悉的星座轮廓将变得难以辨认。$^{[13]}$天文学家可以通过精确测量单个恒星的共同自行$^{[14]}$（依赖高精度天体测量学$^{[15][16]}$）以及通过天文光谱学测定其径向速度$^{[17]}$，来预测星座在过去或未来的形态。

国际天文学联合会所承认的88个星座，以及历史上各文化所使用的星座，都是基于可观测天空中恒星分布而想象出的形象与形状。许多官方承认的星座源自古代近东和地中海地区的神话想象。$^{[18][19][page needed]}$ 其中一些故事似乎与星座的外观有关，例如猎户座被天蝎座刺杀的传说，其对应的星座在一年中的不同季节出现在天空中。$^{[20]}$
\subsection{观测}
\begin{figure}[ht]
\centering
\includegraphics[width=8cm]{./figures/479bf5216e7afbe4.png}
\caption{显示星座、银河以及黄道的星图} \label{fig_Coti_2}
\end{figure}
由于地球围绕太阳运行的轨道上，夜晚发生的时间对应于轨道上的位置不断变化，星座在一年之中其位置也会随之改变。随着地球自转方向指向东方，天球看起来仿佛向西旋转，恒星在北极星周围呈逆时针环绕，在南极星周围则呈顺时针环绕。$^{[21]}$

由于地球自转轴具有23.5°的倾角，黄道沿着黄道面在南、北半球中均匀分布，近似构成一个大圆。北天的黄道星座包括双鱼座、白羊座、金牛座、双子座、巨蟹座和狮子座；南天的黄道星座包括室女座、天秤座、天蝎座、人马座、摩羯座和宝瓶座。$^{[22][a]}$ 根据一年中的不同时间，从北纬23.5°到南纬23.5°的纬度范围内，黄道会出现在头顶正上方。夏季时，黄道在白天显得较高、夜晚较低；而在冬季，情况则正好相反，这一现象在南、北半球均一致。

由于太阳系相对于银河系具有约60°的倾角，银河系的银道面相对于黄道倾斜约60°，位于北天的金牛座与双子座之间，以及南天的天蝎座与人马座之间，银河系中心也位于这一南天区域附近。$^{[22][23]}$ 银河在天空中穿过天鹰座（靠近天赤道），以及北天的天鹅座、仙后座、英仙座、御夫座和猎户座（靠近参宿四），同时也穿过麒麟座（靠近天赤道），以及南天的船尾座、船帆座、船底座、南十字座、半人马座、南三角座和天坛座。$^{[22]}$
\subsubsection{北半球}  
北极星作为北极星，是北天球的近似中心。它属于小熊座，位于小北斗勺柄的末端。$^{[22]}$

在大约北纬35°的地区，一月时，大熊座（包含北斗七星）出现在东北方，而仙后座位于西北方。西方可见双鱼座（位于地平线上方）和白羊座。西南方的鲸鱼座接近地平线。高空偏南的位置可见猎户座和金牛座。东南方向地平线上方是大犬座。在猎户座的上方偏东方向可以看到双子座；东方方向（并逐渐接近地平线）还有巨蟹座和狮子座。除金牛座外，英仙座和御夫座也出现在头顶附近。$^{[22]}$

在相同纬度下，七月时，仙后座（位于低空）和仙王座出现在东北方。大熊座此时位于西北方。牧夫座高悬于西方天空。室女座位于西方，天秤座在西南方，天蝎座在南方。人马座和摩羯座位于东南方。天鹅座（包含北十字星）位于东方。武仙座与北冕座一同高悬于天空中。$^{[22]}$
\begin{figure}[ht]
\centering
\includegraphics[width=6cm]{./figures/58cdd9b51e42621a.png}
\caption{南十字座中的南十字星群以及半人马座中的“南天指针星”可用于寻找南天极星——南极座$\sigma$星。} \label{fig_Coti_3}
\end{figure}
\subsubsection{南半球}  
一月可见的星座包括绘架座和网罟座，分别位于水蛇座和山案座附近。$^{[24]}$

七月时，可以看到天坛座（毗邻南三角座）以及天蝎座。$^{[25]}$

靠近极星的星座包括蝘蜓座、天燕座和南三角座（靠近半人马座），以及孔雀座、水蛇座和山案座。

南极座σ星是最接近南天极的恒星，但在夜空中十分暗弱。因此，可以利用南十字座以及半人马座中的α星和β星（相对于南十字座约逆时针30°）来进行三角定位以确定南天极的位置。$^{[22]}$
\subsection{早期星座的历史}
\subsubsection{法国南部拉斯科洞穴}  
有人提出，位于法国南部拉斯科的、距今约17000年的洞穴壁画描绘了金牛座、猎户座腰带以及昴宿星团等星座。然而，这一观点并未在科学界得到普遍认可。$^{[26][27]}$
\subsubsection{美索不达米亚} 
来自美索不达米亚（今伊拉克地区）的刻石和泥板文书，其年代可追溯至公元前3000年，被认为是人类识别星座的最早、且普遍被接受的证据。$^{[28]}$ 看来，大多数美索不达米亚星座是在相对较短的时间内形成的，大约集中于公元前1300年至公元前1000年之间。美索不达米亚的星座后来出现在许多古典希腊星座体系之中。$^{[29]}$
\subsubsection{古代近东}
\begin{figure}[ht]
\centering
\includegraphics[width=6cm]{./figures/71a76fd0377888ed.png}
\caption{记录公元前164年哈雷彗星的巴比伦泥板} \label{fig_Coti_4}
\end{figure}
最古老的巴比伦恒星与星座目录可追溯至中青铜时代初期，其中尤以《每组三星》文本以及《MUL.APIN》最为著名；后者是在约公元前1000年基于更为精确的观测，对前者加以扩展和修订而成。然而，这些目录中大量出现的苏美尔语名称表明，它们建立在更早但如今已无直接文献留存的早期青铜时代苏美尔传统之上。$^{[30]}$

古典黄道体系是对公元前6世纪新巴比伦星座的修订。希腊人在公元前4世纪采纳了巴比伦的星座体系。托勒密星座中有二十个源自古代近东，另有十个使用相同的恒星，但名称不同。$^{[29]}$

《圣经》学者 E.\,W.\,Bullinger 将《以西结书》和《启示录》中提及的一些生物解释为黄道四分区的中间星座标志$^{[31][32]}$：狮子对应狮子座，公牛对应金牛座，人象征宝瓶座，而鹰则代表天蝎座。$^{[33]}$ 《圣经·约伯记》中也提及若干星座，包括 עיש ‘Ayish（意为“灵柩”）、כסיל chesil（“愚者”）以及 כימה chimah（“堆积”）（约伯记 9:9，38:31–32），在《钦定版圣经》中被译为“大角星、猎户座和昴宿星团”，但其中 ‘Ayish“灵柩” 实际对应的是大熊座。$^{[34]}$ 术语 Mazzaroth מַזָּרוֹת，被译为“冠冕的花环”，在《约伯记》38:32中是一个仅出现一次的词，其含义可能指黄道星座。
\subsubsection{古典古代}
\begin{figure}[ht]
\centering
\includegraphics[width=6cm]{./figures/8bd69f4efca5ef5d.png}
\caption{埃及星图与旬星时钟，出自塞嫩穆特墓穴天花板，约公元前1473年} \label{fig_Coti_5}
\end{figure}
关于古希腊星座的信息所存不多，仅有一些零散证据可见于希腊诗人赫西俄德的《工作与时日》，其中提到了“天体”。$^{[35]}$在希腊化时期，希腊天文学基本上采纳了更早的巴比伦体系，该体系最早由克尼多斯的欧多克索斯于公元前4世纪引入希腊。$^{[citation needed]}$ 欧多克索斯的原著已佚，但其内容以诗体形式保存在公元前3世纪阿拉托斯的作品中。现存关于星座神话起源最为完整的著作，出自一位被称为“伪埃拉托色尼”的希腊化时期作家，以及一位被称为“伪许吉努斯”的早期罗马作家。从古代晚期直到近代早期，西方天文学教学的基础是托勒密于公元2世纪撰写的《天文学大成》。

在托勒密王国时期，埃及本土传统以拟人化形象来表现行星、恒星及各种星座。$^{[36]}$ 这些传统后来与希腊和巴比伦的天文学体系相结合，最终形成了丹德拉黄道，这是目前已知最早完整描绘全部现今熟知黄道星座的实例，同时还包含一些原创的埃及星座、旬星以及行星。$^{[28][37]}$ 这一融合发生的具体时间仍不明确，但多数观点认为其形成于罗马时期，即公元2至4世纪之间。托勒密的《天文学大成》在中世纪时期仍是欧洲及伊斯兰天文学中界定星座的标准依据。$^{[38]}$
\subsubsection{中国古代}
\begin{figure}[ht]
\centering
\includegraphics[width=8cm]{./figures/4fe77cf6fcfe09c6.png}
\caption{采用圆柱投影的中国星图（苏颂）} \label{fig_Coti_6}
\end{figure}
中国古代具有悠久的天象观测传统。$^{[39]}$ 一些并不具体的中国恒星名称，后来被归入二十八宿体系，已在出土于安阳、可追溯至商代中期的甲骨文中发现。中国星座体系是对中国天空最为重要的观测成果之一，其历史可追溯至公元前5世纪。中国的星座体系独立于希腊—罗马体系而发展，尽管通过与古巴比伦天文学的若干相似之处，推测二者之间可能存在较早的相互影响。$^{[40]}$

汉代经典中国天文学的三大学派，被认为源自更早的战国时期天文学家。三大学派的星座体系在三国时期、3世纪的天文学家陈卓手中被整合为一个统一体系。陈卓的原著已佚，但其星座体系的信息保存在唐代文献中，尤以居住在长安的中印天文学家瞿昙悉达所著《九执历》（Chinese: 九執曆；pinyin: Jiǔzhí-lì）为代表。$^{[41]}$ 现存最早的中国星图亦可追溯至这一时期，并作为敦煌文献的一部分保存至今。中国本土天文学在宋代达到繁荣，而在元代以及随后的明代，则逐渐受到中世纪时期其他东方和西方天文学家的影响，参见《开元占经》。$^{[40]}$ 由于这一时期绘制的星图采用了更具科学性的方式，因此被认为更为可靠。$^{[42]}$

宋代一幅著名的星图是《苏州天文图》，该图将中国天空的恒星雕刻在一块石制平面天球图上，基于实际观测精确绘制，并标示了位于金牛座的公元1054年超新星。$^{[42]}$

明代晚期在欧洲天文学的影响下，星图中描绘了更多恒星，但仍保留了传统星座体系。新观测到的恒星被作为补充并入南天的旧有星座之中，这些南天区域并未出现在古代中国天文学家的传统记录中。明代后期，徐光启与德国耶稣会士汤若望（Johann Adam Schall von Bell）进一步改进了这一体系，其成果记录于《崇祯历书》（1628年）。$^{[clarification needed]}$ 传统中国星图基于西方星图的知识，引入了南天23个新星座、共125颗恒星；随着这一改进，中国的天空体系得以与世界天文学接轨。$^{[42][43]}$
\subsection{近代早期天文学}  
从历史角度看，北天与南天星座的起源具有显著差异。大多数北天星座可追溯至古代，其名称主要源自古典希腊神话。$^{[11]}$ 关于这些星座的证据以星图形式保存下来，其中最早的表现形式可能出现在被称为“法尔内塞天球仪”的雕像上，该雕像或许基于希腊天文学家喜帕恰斯的星表。$^{[44]}$ 南天星座则多为近代的创造，有时作为古代星座的替代（例如阿尔戈船座）。一些南天星座最初名称较长，后来被缩短为更易使用的形式，例如 Musca Australis 最终简化为 Musca。$^{[11]}$

一些早期提出的星座从未被普遍接受。不同观测者常以不同方式将恒星组合成星座，而任意划定的星座边界也常导致对天体归属星座的混淆。在天文学家开始精确划定星座边界之前（始于19世纪），星座通常只是天空中界限模糊的区域。$^{[45]}$ 如今，星座边界遵循官方认可的赤经与赤纬分界线，这些界线基于本杰明·古尔德在其星表《阿根廷天区测量》（Uranometria Argentina）中所定义的 B1875.0 春分点。$^{[46]}$

1603年，约翰·拜耳出版的星图集《天图》（Uranometria）将恒星明确归属至各个星座，并通过为每个星座内的恒星分配希腊字母与拉丁字母，正式确立了这一划分体系。这种命名方式如今被称为拜耳命名法。$^{[47]}$ 随后的星图编制进一步促成了当今通行的现代星座体系。

若干近代提出的星座方案并未延续下来。例如，法国天文学家皮埃尔·勒莫尼耶和约瑟夫·拉朗德曾提出一些一度流行、但后来被废弃的星座。北天星座“象限仪座”（Quadrans Muralis）曾沿用至19世纪，其名称后来被用于象限仪座流星雨，但如今该星座已被划分并并入牧夫座与天龙座之中。$^{[48]}$
\subsubsection{南天星座的起源}
\begin{figure}[ht]
\centering
\includegraphics[width=8cm]{./figures/8f18b38336c2466a.png}
\caption{} \label{fig_Coti_7}
\end{figure}
